\subsection{Schism: workload-aware разбиение графовых баз данных}

Schism \cite{schism} представляет собой метод разбиения графовых и реляционных структур, ориентированный на минимизацию стоимости выполнения рабочей нагрузки запросов. В отличие от классических алгоритмов partitioning, которые опираются исключительно на структуру графа, Schism использует информацию о реальных запросах для построения модели взаимодействий между данными. Если быть конкретнее, то основной целью Schism является уменьшение числа межузловых обращений, возникающих при обработке рабочей нагрузки.

\subsubsection{Постановка задачи Schism}

Рассматривается граф данных $G = (V, E)$, где $V$ представляет множество вершин, а $E$ множество рёбер. Также задано множество запросов $Q = \{q_1, q_2, \dots, q_m\}$, отражающее рабочую нагрузку системы.

Каждый запрос $q_i \in Q$ представляет собой множество обращений к вершинам графа, которые могут выполняться последовательно или в рамках одного транзакционного контекста.

Требуется построить разбиение множества вершин $V$ на $k$ непересекающихся подмножеств $P_1, P_2, \dots, P_k$:

\begin{equation}
\bigcup_{i=1}^{k} P_i = V, \quad P_i \cap P_j = \emptyset, \quad i \neq j,
\end{equation}

такое, что минимизируется стоимость выполнения рабочей нагрузки:

\begin{equation}
C(Q) = \sum_{q \in Q} C(q) \rightarrow \min.
\end{equation}

При этом стоимость отдельного запроса определяется количеством межпартиционных обращений, возникающих при его выполнении.

Дополнительно вводится ограничение балансировки размеров партиций:

\begin{equation}
|P_i| \approx \frac{|V|}{k}.
\end{equation}

\textbf{Общая идея Schism}

Основная идея Schism заключается в переходе от структурного разбиения графа к разбиению, основанному на анализе рабочей нагрузки. Вместо использования только структуры исходного графа, метод использует статистику запросов для построения модели совместного использования данных.

Интуитивно, если две вершины часто участвуют в одних и тех же запросах, их размещение в одной партиции позволяет снизить количество распределённых операций. Таким образом, Schism стремится минимизировать число запросов, пересекающих границы партиций.

\subsubsection{Модель рабочей нагрузки}

Для построения модели рабочей нагрузки используется журнал запросов $Q$. На основе этого журнала формируется вспомогательная структура, отражающая частоту совместного доступа к вершинам графа.

Для пары вершин $u, v \in V$ вводится функция совместного использования:

\begin{equation}
w(u, v) = |\{q \in Q \mid u \in q \land v \in q\}|.
\end{equation}

Данная функция определяет количество запросов, в которых одновременно участвуют вершины $u$ и $v$. Чем выше значение $w(u, v)$, тем более выгодно размещать данные вершины в одной партиции.

На основе данной функции строится взвешенный граф рабочей нагрузки:

\begin{equation}
G_w = (V, E_w),
\end{equation}

где $E_w$ включает пары вершин с ненулевым значением $w(u, v)$.

\textbf{Интерпретация графа рабочей нагрузки}

Граф $G_w$ не отражает структуру исходного графа данных, а моделирует поведение пользователей системы. Вершины данного графа соответствуют сущностям исходной базы данных, а рёбра отражают частоту их совместного использования в рамках запросов.

Вес ребра интерпретируется как мера зависимости между данными объектами с точки зрения рабочей нагрузки. Таким образом, задача разбиения сводится к минимизации разрезов именно в этом графе, а не в исходной структуре $G$.

\textbf{Построение модели оптимизации}

Schism преобразует задачу разбиения в задачу оптимизации, где целевая функция соответствует стоимости выполнения рабочей нагрузки. Для каждого запроса $q \in Q$ определяется количество межпартиционных обращений, возникающих при его выполнении.

Пусть запрос $q$ затрагивает множество вершин, распределённых по различным партициям. Тогда стоимость запроса определяется как функция количества переходов между партициями:

\begin{equation}
C(q) = f(\{P_i \mid v \in q, v \in P_i\}),
\end{equation}

где функция $f$ отражает количество межузловых взаимодействий.

\textbf{Инициализация разбиения}

На начальном этапе алгоритма формируется начальное разбиение $\Pi^{(0)} = \{P_1^{(0)}, \dots, P_k^{(0)}\}$. Данное разбиение может быть получено случайным образом либо с использованием простой эвристики, основанной на локальных свойствах графа.

Качество начального разбиения не является критическим фактором, поскольку дальнейшая оптимизация выполняется итеративно.

\textbf{Итеративная оптимизация}

Основной этап Schism заключается в итеративной локальной оптимизации разбиения. Для каждой вершины $v \in V$ рассматривается возможность её перемещения из текущей партиции $P_i$ в другую партицию $P_j$.

Изменение принимается, если выполняется условие улучшения:

\begin{equation}
\Delta C < 0,
\end{equation}

где $\Delta C$ обозначает изменение стоимости рабочей нагрузки после перемещения вершины.

Процесс повторяется до достижения состояния, в котором дальнейшие улучшения становятся невозможными или незначительными.

\textbf{Балансировка партиций}

Помимо минимизации стоимости запросов, алгоритм Schism учитывает ограничение на размер партиций:

\begin{equation}
|P_i| \leq (1 + \epsilon)\frac{|V|}{k}.
\end{equation}

Это ограничение предотвращает возникновение дисбаланса между узлами хранения, который может привести к снижению производительности системы.

При нарушении ограничения выполняются корректирующие перемещения вершин между партициями.

\subsubsection{Вычислительная сложность}

Основная вычислительная сложность Schism связана с построением графа рабочей нагрузки и оценкой стоимости запросов. При наивной реализации построение $G_w$ требует анализа всех пар вершин, участвующих в запросах, что может приводить к высокой вычислительной нагрузке.

Итеративная оптимизация имеет сложность, зависящую от числа вершин $|V|$, числа партиций $k$ и числа итераций $I$:

\begin{equation}
O(I \cdot |V| \cdot k).
\end{equation}

\textbf{Свойства получаемого разбиения}

Получаемое в результате работы Schism разбиение обладает свойством локализации рабочей нагрузки, при котором вершины, часто используемые совместно, стремятся находиться в одной партиции.

Это приводит к снижению числа межузловых обращений при выполнении запросов и повышению эффективности обработки рабочей нагрузки в распределённых системах.

\subsubsection{Ограничения метода}

Несмотря на эффективность, Schism имеет ряд ограничений. Основным из них является высокая стоимость построения точной модели рабочей нагрузки, требующей анализа большого количества пар взаимодействий между вершинами.

Кроме того, метод предполагает относительно стабильную рабочую нагрузку, так как существенные изменения характера запросов требуют пересчёта разбиения.

Также следует отметить, что при увеличении размера графа и числа запросов возрастает как вычислительная, так и память-затратная сложность построения модели взаимодействий.

\subsubsection{Сравнение с классическими методами}

В отличие от классических методов edge-cut разбиения, которые минимизируют число разрезанных рёбер исходного графа, Schism минимизирует стоимость выполнения рабочей нагрузки.

Таким образом, если традиционные методы оптимизируют выражение

\begin{equation}
\min |E_{\text{cut}}|,
\end{equation}

то Schism оптимизирует выражение

\begin{equation}
\min C(Q).
\end{equation}

Это делает Schism более чувствительным к реальному поведению системы, но одновременно увеличивает сложность построения модели разбиения.